\label{subsec:xi-certificate-coherence-orthogonal-axis-capacity-part1}

This subsection continues the structural reduction concerning "endpoint quadratic compression/bit-budget barrier—certificate coherence—falsifiable interface" described in the overview subsection \ref{subsubsec:xi-ramanujan-horizon-ledger-framework}: building upon the three theorem chains—defect exclusion radius, timefiber endpoint compression, and \(\mathbf{PW}\) no-continuous-hair—established in subsection \ref{subsec:dephysicalized-horizon-timefiber-nohair}, and under the umbrella of the analytic—calculus—arithmetic/spectral triple-certificate closed-loop formulation of Theorem \ref{thm:dephys-rh-certificate-audit-loop}, we further extract propositions directly aligned with "verifiable semantics." The core is to rigorize the zero-events on the unitary slice as an inverse system of \emph{finite certificate objects}, and to compress the information "off the slice" into the endpoint second-order radial channel and the capacity lower bound induced by the precision potential.

\paragraph{Three-axis addressing and verifiable semantics}
The readout in this paper follows the convention "address precedes value," and compiles non-verifiable objects into \(\mathsf{NULL}\) via visible-domain constraints and checkers (Definition \ref{def:xi-visible-read-null}).
Under this convention, addresses are not auxiliary annotations but rather \emph{internalized readable prefixes}: if an address object is not given, then there exists no evaluation function \(\ind_{C_a}\) referenceable at the object layer (equivalently, the evaluation function \(1_{C_a}\) does not hold at the type level), and the readout partial function is undefined at the type level and represented as \(\mathsf{NULL}\) (\S\ref{sec:recursive-addressing}; see also the \(\mathscr{F}(a)=\varnothing\) calibration of Proposition \ref{prop:null-as-local-section-obstruction}).
At the same time, recursive addressing does not expand the existing visible \(\sigma\)-algebra: the derived cylinder event family falls strictly within the old-layer visible domain \(\mathcal{G}^{(L)}\), merely re-selecting and organizing existing observable events, so that \(\mathcal{G}^{(L+1)}\subseteq \mathcal{G}^{(L)}\) (Proposition \ref{prop:recursive-addressing-refinement} and Corollary \ref{cor:recursive-addressing-visible-sigma-nonexpansion}).
Meanwhile, a verifiable address can be understood as a multi-axis prefix
\(
a=(a_{\mathrm{vis}},a_{\mathrm{res}},a_{\mathrm{mode}})
\)
(Definition \ref{def:multi_axis_address}):
where \(a_{\mathrm{vis}}\) aligns with the single-parameter phase interface of the unitary slice (Proposition \ref{prop:xi-unitary-slice-spectral-crossing}),
\(a_{\mathrm{res}}\) carries the residue/inverse-limit axis externalization needed to avoid image collisions (Proposition \ref{prop:xi-null-three-way} and Theorem \ref{thm:address-defect-law}),
and \(a_{\mathrm{mode}}\) records the orthogonal mode axes allowed by the protocol (typically including radial registers and minimal realization dimension axes; see Corollary \ref{cor:xi-null-second-order-radial-channel} and Remark \ref{rem:xi-acc-rank}).
The conclusions below correspond in turn to: the coherent merger of certificates on the \(a_{\mathrm{vis}}\) axis, the type consequences of the unitary slice as an orthogonal slice, the endpoint capacity threshold for off-slice information, and the resulting falsifiable interface restatement.

\begin{remark}[Terminology convention: axis/channel versus projective system automorphism]\label{rem:axis-vs-automorphism-terminology}
In this paper, the term "axis" is used solely as a shorthand for \emph{addressable coordinates/registers} (and their budget directions) in protocol semantics, referring to different record components in address prefixes (e.g., \(a_{\mathrm{vis}},a_{\mathrm{res}},a_{\mathrm{mode}}\)), rather than as symmetry objects.
The symmetry objects discussed in this paper are the \emph{automorphisms} of the projective readout system and their descent actions on visible sections; in the gauge groupoid calibration, these can be understood as the actions generated by composable morphisms that preserve projective gates and verification interfaces.
Therefore, "orthogonal decomposition/channels" appearing in the text all refer to invariant subrepresentations, direct sum decompositions, or symbolic decompositions under the automorphism action: for example, the modulo-\(4\) grouping of endpoint payload and the decoupling of precision potential ledger (Proposition \ref{prop:endpoint-orthogonal-decomposition-z4-Rplus}), as well as the one-dimensional eigensubrepresentation of endpoint contravariant observables and their \(L(1+\sigma,\chi_4)\) symbol (Theorem \ref{thm:endpoint-zeta-Lchi4-symbol-pair}).
\end{remark}

\begin{theorem}[Two-symbol jet of endpoint regularized power chain: \texorpdfstring{$\zeta$}{zeta}-type tristate constant layer \(\perp\) \texorpdfstring{$\chi_4$}{chi4}-type linear hair layer]\label{thm:xi-endpoint-two-symbol-jet}
Let \(g\) be analytic on an open neighborhood containing the interval \([-2,2]\).
Take the completed trace coordinate \(S\) and Chebyshev--Adams polynomials \(C_m(S)\) (Proposition \ref{prop:half-angle-z4-residue}), and take \(\sigma>1\).
Following the endpoint regularized power chain operator \(\mathcal W_\sigma\) in the appendix (equation \eqref{eq:endpoint-Wsigma-even}), and naturally extending it to locally analytic test functions \(g\):
\[
(\mathcal W_\sigma g)(S):=\sum_{m\ge 1}\frac{1}{m^{1+\sigma}}\,g\!\bigl(C_m(S)\bigr).
\]
Then as \(S\to 0\), the endpoint jet expansion holds:
\begin{equation}\label{eq:xi-endpoint-two-symbol-jet}
\boxed{\;
(\mathcal W_\sigma g)(S)
:=A_\sigma(g)\;+\;L(\sigma,\chi_4)\,g'(0)\,S\;+\;O(S^2)\;}
\end{equation}
where the constant layer depends only on the endpoint tristate values \(g(0),g(2),g(-2)\):
\[
A_\sigma(g):=\mathsf H_{\mathrm{odd}}(\sigma)\,g(0)+\mathsf H_{0}(\sigma)\,g(2)+\mathsf H_{2}(\sigma)\,g(-2),
\]
and the three-channel weight kernels are
\[
\mathsf H_{\mathrm{odd}}(\sigma):=\sum_{\substack{m\ge 1\\ m\ \text{odd}}}\frac{1}{m^{1+\sigma}},\qquad
\mathsf H_{0}(\sigma):=\sum_{m\equiv 0(4)}\frac{1}{m^{1+\sigma}},\qquad
\mathsf H_{2}(\sigma):=\sum_{m\equiv 2(4)}\frac{1}{m^{1+\sigma}}
\]
(see Proposition \ref{prop:endpoint-dirichlet-symbol-weights} for their closed-form \(\zeta\)-symbolic representation).
The first-order jet coefficient is controlled by the modulo-\(4\) character \(\chi_4\): let
\[
\chi_4(m):=\sin(m\pi/2)\in\{-1,0,1\},
\qquad
L(\sigma,\chi_4):=\sum_{m\ge 1}\frac{\chi_4(m)}{m^\sigma},
\]
then \(\chi_4\) is equivalent to the nontrivial Dirichlet character modulo \(4\) (Theorem \ref{thm:endpoint-four-channel-chi4}), and the linear term is fully induced by the first-order jet \(C_m'(0)=m\sin(m\pi/2)\) of the odd channel.
This first-order jet introduces an additional factor \(m\), thereby shifting the weight exponent \(1+\sigma\) of \(\mathcal W_\sigma\) to \(L(\sigma,\chi_4)\) at the symbolic level.
\end{theorem}

\begin{proof}
By the endpoint tristate compression \(C_m(0)\in\{0,\pm 2\}\) (Proposition \ref{prop:half-angle-z4-residue} and Corollary \ref{cor:endpoint-tristate-compression}), and by the odd-even jet expansion of Corollary \ref{cor:chebyshev-endpoint-jet-regularization},
when \(S\to 0\),
\[
g(C_m(S))=
\begin{cases}
g(C_m(0))+O(S^2), & m\ \text{even},\\
g(0)+g'(0)\,C_m'(0)\,S+O(S^2), & m\ \text{odd}.
\end{cases}
\]
For the even channels grouped by \(m\bmod 4\), the constant layer is a linear combination of the tristate weight kernels; Proposition \ref{prop:endpoint-dirichlet-symbol-weights} gives their closed-form \(\zeta\)-symbolic representation.
For the linear term, even \(m\) contribute no first-order term, while odd \(m\) yield by Proposition \ref{prop:chebyshev-endpoint-derivative-splitting}
\(
C_m'(0)=m\sin(m\pi/2)=m\chi_4(m)
\),
so the linear coefficient is
\[
\sum_{m\ge 1}\frac{1}{m^{1+\sigma}}\,C_m'(0)
:=\sum_{m\ge 1}\frac{\chi_4(m)}{m^\sigma}
:=L(\sigma,\chi_4).
\]
Combining the above yields \eqref{eq:xi-endpoint-two-symbol-jet}. \qed
\end{proof}

\begin{corollary}[Endpoint orthogonal decoupling (jet level): logarithmic main-scale constant layer \(\perp\) \texorpdfstring{$\chi_4$}{chi4} linear memory]\label{cor:xi-endpoint-jet-orthogonal-decoupling}
Under the setting of Theorem \ref{thm:xi-endpoint-two-symbol-jet}, the two types of information in the endpoint jet are strictly separated:
\begin{enumerate}
  \item The constant layer \(A_\sigma(g)\) is determined solely by the modulo-\(4\) grouping weight kernels of endpoint tristates \(C_m(0)\in\{0,\pm 2\}\), thus belonging to the \(\zeta\)-type main-scale channel (Proposition \ref{prop:endpoint-dirichlet-symbol-weights}).
  \item The first-order jet is fully controlled by the \(\chi_4\) channel \(L(\sigma,\chi_4)\), and does not couple with the additive shift law of the endpoint precision potential ledger \(y\) (Proposition \ref{prop:endpoint-orthogonal-decomposition-z4-Rplus}).
\end{enumerate}
Therefore, at the endpoint jet level, "the continuous infinity logarithmic ledger" and "the linear memory of discrete branching hairs" exhibit strict orthogonal decoupling: the former enters only the \(\zeta\)-symbolic weight kernels of the tristate constant layer, while the latter occupies the unique first-order visible branch in the form of a one-dimensional \(\chi_4\) eigenchannel.
\end{corollary}

\begin{proof}
By Theorem \ref{thm:xi-endpoint-two-symbol-jet}, the constant layer depends only on endpoint tristate values and its weight kernels are given by \(\zeta(1+\sigma)\)-type Dirichlet regularization; while the linear layer is given solely by the \(\chi_4\) channel.
Proposition \ref{prop:endpoint-orthogonal-decomposition-z4-Rplus} shows the orthogonal decomposition "modulo-\(4\) discrete payload \(\oplus\) precision potential ledger" at the endpoint, hence the two jet components do not couple. \qed
\end{proof}

\begin{theorem}[Quadratic resolution barrier of Jensen defect certificate: unavoidable radius threshold for finite-height bands]\label{thm:xi-jensen-quadratic-resolution-barrier}
Fix \(T\ge 1\) and \(\delta_0\in(0,\tfrac12]\), and denote
\[
R(T,\delta_0)^2:=\frac{T^2+(\delta_0-1)^2}{T^2+(\delta_0+1)^2}\in(0,1)
\]
(Theorem \ref{thm:xi-jensen-single-radius-band-exclusion}).
If a certain RH certificate chain takes only the defect upper bound chain \(D_\zeta(\varrho_n)\le \varepsilon_n\to 0\) (equivalently \(r_\ast(\varrho_n)=\varrho_n e^{-D_\zeta(\varrho_n)}\uparrow 1\), see Theorem \ref{thm:dephys-defect-repulsion-radius}) as the sole input,
and aims to exclude all off-slice zeros satisfying offset lower bound \(\tfrac12-\beta\ge\delta_0\) in the height band \(|\gamma|\le T\),
then asymptotically its radius sequence must enter the unavoidable quadratic resolution layer:
\begin{equation}\label{eq:xi-jensen-quadratic-resolution-barrier}
\boxed{\;
1-\varrho_n\ \lesssim\ \frac{\delta_0}{T^2}\;}
\end{equation}
(constant depends only on boundedness constraints of \(\delta_0\)).
\end{theorem}

\begin{proof}
By Theorem \ref{thm:xi-offcritical-quadratic-radial-compression},
any off-slice zero \(\rho=\beta+\mathrm{i}\gamma\) satisfying \(|\gamma|\le T\) and \(\delta=\tfrac12-\beta\ge\delta_0\) has disk image \(w_\rho\) satisfying
\(|w_\rho|\le R(T,\delta_0)\);
therefore if \(\varrho\le R(T,\delta_0)\), then such off-line zeros do not necessarily enter the "in-disk contribution" term of the Jensen summation (Proposition \ref{prop:app-jensen-formula-defect}), so the upper bound on \(D_\zeta(\varrho)\) alone cannot form a band-exclusion conclusion.
To achieve structural exclusion within this height band, the certificate chain must reach \(\varrho_n>R(T,\delta_0)\) at some level (certificate form of Theorem \ref{thm:xi-jensen-single-radius-band-exclusion}).
\par
On the other hand, direct calculation gives
\[
1-R(T,\delta_0)^2
:=\frac{4\delta_0}{T^2+(\delta_0+1)^2}
\asymp \frac{\delta_0}{T^2},
\]
so once \(\varrho_n>R(T,\delta_0)\), we must have
\(
1-\varrho_n\le 1-R(T,\delta_0)\lesssim 1-R(T,\delta_0)^2\lesssim \delta_0/T^2
\),
yielding \eqref{eq:xi-jensen-quadratic-resolution-barrier}. \qed
\end{proof}

\begin{corollary}[Irreplaceability of endpoint angle window: high-height contribution forces localization at \texorpdfstring{$\theta\approx\pi$}{theta≈pi}]\label{cor:xi-jensen-endpoint-angle-window-rigidity}
At the resolution layer \(1-\varrho\asymp \delta_0/T^2\) of Theorem \ref{thm:xi-jensen-quadratic-resolution-barrier},
the high-imaginary-part contribution of the Jensen circle average is forced to concentrate in the endpoint angle window \(\theta\approx \pi\) (i.e., \(w=\varrho e^{\mathrm{i}\theta}\approx -1\)):
the preimage of the equal-radius fiber in the upper half-plane simultaneously touches \(y\downarrow 0\) and \(y\uparrow \infty\), and the maximum height \(y_+(\varrho)=\frac{1+\varrho}{1-\varrho}\) is attained only at \(\theta=\pi\) (Theorem \ref{thm:dephys-timefiber-two-ended}).
Therefore, any effective band information based on defect circle average/certificate upper bound chains must enter the closed-loop interface through the endpoint angle window and endpoint jet (Theorem \ref{thm:xi-endpoint-two-symbol-jet}), and cannot be equivalently replaced by uniform angular sampling away from the endpoint.
\end{corollary}

\begin{proof}
By Theorem \ref{thm:dephys-timefiber-two-ended}, the imaginary part of the Cayley preimage is
\(
\Im t(\varrho,\theta)=\frac{1-\varrho^2}{1+\varrho^2+2\varrho\cos\theta}
\),
and attains its maximum at \(\theta=\pi\) as \(y_+(\varrho)=\frac{1+\varrho}{1-\varrho}\to\infty\) (\(\varrho\uparrow 1\)).
Therefore, when \(\varrho\) enters the endpoint resolution layer, any main contribution involving height windows must come from the neighborhood \(\theta\approx\pi\); combining this endpoint localization with endpoint jet symbolic decomposition yields the conclusion. \qed
\end{proof}

\begin{theorem}[Triple consistency audit principle: \texorpdfstring{$p$}{p}-adic closure \(\times\) Hankel rigidity \(\times\) semantic conservativity]\label{thm:xi-triple-consistency-audit}
Under the auditable semantics calibration of bounded slice \(\mathbf{PW}^{\ast}_{\le}\) (Definition \ref{def:pom-sem-functor}),
any implementation claiming to output "spectral fingerprint (such as \(r_q(u)\) or Hankel-normal form)—\(\zeta\)-closed-loop—defect certificate chain (such as \(r_\ast(\varrho_n)\uparrow 1\))"
must, to be regarded as verifiable within the system, simultaneously satisfy three mutually orthogonal verification constraints:
\begin{enumerate}
  \item \textbf{\((p)\)-adic closure guardrails:} Integer-layer constraints such as Dwork congruence/sparsification/endpoint vanishing and \(p=2\) Chebyshev--Dwork single-variable closed recursion (Proposition \ref{prop:dephys-padic-guardrails-endpoint-rigidity}).
  \item \textbf{Hankel minimal realization rigidity:} Minimal offset certificate and \(2d\)-point locking Hankel-normal form, thereby fixing minimal realization dimension and spectral radius fingerprint as finite computable objects (Theorem \ref{thm:dephys-hankel-finite-audit}; see also Remark \ref{rem:xi-acc-rank} for dimension-type \(\mathsf{NULL}\) interpretation).
  \item \textbf{Semantic conservativity and finite signature compression:} The conservative semantic functor \(\mathsf{Sem}\) constrains implementation differences to invertible \(2\)-morphisms, so that the above fingerprints and guardrails possess implementation-independence in signature/spectral domains (Theorem \ref{thm:pom-sem-conservative}).
\end{enumerate}
Thus, "information does not leak" under this calibration is no longer an empirical assertion, but a computable, falsifiable constraint system jointly enforced by three mutually supporting guardrails (see Theorem \ref{thm:dephys-rh-certificate-audit-loop} for the analytic—calculus—arithmetic/spectral triple-certificate closed loop).
\end{theorem}

\begin{proof}
Conditions (1) and (2) are given by Proposition \ref{prop:dephys-padic-guardrails-endpoint-rigidity} and Theorem \ref{thm:dephys-hankel-finite-audit}, respectively;
condition (3) is given by Theorem \ref{thm:pom-sem-conservative} (slice conservativity).
The orthogonality of the three constraints comes from their supporting output domains being mutually non-overlapping: \((p)\)-adic guardrails occur at the integer congruence layer, Hankel rigidity occurs at the finite-window/minimal-realization layer, and semantic conservativity occurs at the natural isomorphism class layer of slices; their merger yields the verifiable closed-loop audit principle. \qed
\end{proof}

\begin{remark}[Closed-loop reduction (verifiable formulation)]\label{rem:xi-auditable-reduction-closed-loop}
Theorems \ref{thm:xi-endpoint-two-symbol-jet}--\ref{thm:xi-triple-consistency-audit} provide a closed-loop reduction:
the endpoint power chain splits at the jet level into a \(\zeta\)-type tristate constant layer and a \(\chi_4\)-type linear hair layer (Corollary \ref{cor:xi-endpoint-jet-orthogonal-decoupling}),
the defect certificate chain has an unavoidable quadratic resolution barrier and forces endpoint angle window main rigidity (Theorem \ref{thm:xi-jensen-quadratic-resolution-barrier} and Corollary \ref{cor:xi-jensen-endpoint-angle-window-rigidity}),
and any implementation to carry "RH certificatization" must pass the triple consistency audit of \((p)\)-adic closure, Hankel rigidity, and semantic conservativity (Theorem \ref{thm:xi-triple-consistency-audit}).
These three points jointly elevate the claim "off-slice information can only be orthogonally externalized" to a computable, falsifiable structural constraint, and are strictly compatible with the semantic premise "address precedes value/recursive visible domain does not expand" at the beginning of this section (\S\ref{sec:recursive-addressing}).
\end{remark}

\paragraph{Extension-preserving category: canonical record axes and unique continuous transverse direction}
The "extension-preserving" language used in this section aims to view "irreversible readout" as a structural phenomenon of \emph{insufficient record axes}, and to rigorize information conservation as injectivity on the extended space by externalizing record axes (Proposition \ref{prop:unit-circle-phase-extension-preserving}; see also Appendix \ref{subsec:unit-circle-phase-information}).
Under verifiable semantics, this idea can be organized as a quotient category: objects are record axes that make extended readout injective, and morphisms are compressions/projections that preserve certificate-verifiable equivalence.

\begin{definition}[Extension-preserving category and quotient morphisms]\label{def:xi-extension-preserving-category}
Fix a readout (allowing many-to-one compression) \(\pi:\Omega\to Y\), and take "verifiable single-valued readout" as a type-level premise (Proposition \ref{prop:xi-null-three-way}).
Call a map \(r:\Omega\to R\) a \textbf{record axis}, and call the extended readout
\[
\pi^{\mathrm e}:=(\pi,r):\Omega\to Y\times R
\]
\textbf{extension-preserving} if \(\pi^{\mathrm e}\) is injective at the object level (equivalently: \(\omega\neq\omega'\Rightarrow (\pi(\omega),r(\omega))\neq(\pi(\omega'),r(\omega'))\)).
\par
Define \(\mathrm{Ext}(\pi)\) as the following category: objects are record axes \(r\) that make \(\pi^{\mathrm e}\) injective;
morphisms \(r\to r'\) are maps \(\phi:R\to R'\) such that \(r'=\phi\circ r\), and require this factorization to preserve coherence in the certificate-verifiable equivalence sense (i.e., \(\phi\) does not introduce new non-offline-verifiable distinctions; in concrete implementations this can be taken as quotient maps "encoded by finite objects and reproducible by deterministic checkers").
\end{definition}

\begin{theorem}[Initial object of extension-preserving category and canonical record axis]\label{thm:xi-extension-preserving-initial-object}
Under unitary slice visible domain constraints and \(\mathbf{PW}\) closure calibration, further assume:
\begin{enumerate}
  \item \textbf{(Orthogonal trichotomy)} The sources of global \(\mathsf{NULL}\) decompose within the system into mutually non-substitutable orthogonal witnesses (Theorem \ref{thm:xi-null-orthogonal-trichotomy-resource-nonsubstitutable});
  \item \textbf{(Rigidity of minimal externalization)} Reversibility of image collision requires externalizing residue/inverse-limit axis and satisfying defect budget lower bound (Theorem \ref{thm:address-defect-law}), reversibility of imaginary folding requires externalizing winding certificate (Theorem \ref{thm:dephys-dirichlet-winding-minimal}), and continuous accumulation of endpoint contraction semigroup can occur only along the unique Koenigs depth ledger (Theorem \ref{thm:dephys-endpoint-koenigs-rigidity});
  \item \textbf{(Closed-form sealing of endpoint finite part)} All divergences of endpoint even channels are compressed after renormalization into universal \(\zeta\)-pole and explicit finite part, and are orthogonal to the contravariant \(\chi_4\) eigenchannel (Theorem \ref{thm:endpoint-even-finite-part-renormalization} and Remark \ref{rem:endpoint-renormalization-universal-pole-discrete-hair}).
\end{enumerate}
Then in the extension-preserving category \(\mathrm{Ext}(\pi)\) (Definition \ref{def:xi-extension-preserving-category}) there exists a canonical object
\[
r_0
:=
r_{\mathrm{fin}}\times r_{\widehat{\ZZ}}\times r_{\mathrm{wind}}\times r_{\mathrm{end}},
\]
where \(r_{\mathrm{fin}}\) records the universal pole/finite part of endpoint even channels, \(r_{\widehat{\ZZ}}\) is the residue/inverse-limit (prime-register) record axis (finite-support truncation calibration of congruence axis \(\widehat{\ZZ}\) see Remark \ref{rem:pom-mod-support-count}), \(r_{\mathrm{wind}}\) is the winding certificate axis (\(\kappa_\lambda\), equation \eqref{eq:dirichlet-winding-certificate}), and \(r_{\mathrm{end}}\) is the endpoint Koenigs depth ledger (Theorem \ref{thm:dephys-endpoint-koenigs-rigidity}).
Moreover, for any object \(r\in\mathrm{Ext}(\pi)\), there exists a unique (in certificate-verifiable equivalence sense) morphism \(r_0\to r\).
In other words: \(r_0\) is the initial object of \(\mathrm{Ext}(\pi)\), and any other extension-preserving record axis can be viewed as a factorization/compression of \(r_0\).
\end{theorem}

\begin{proof}
By Theorem \ref{thm:xi-null-orthogonal-trichotomy-resource-nonsubstitutable}, failure witnesses under orthogonal decomposition are mutually non-substitutable, so any extension-preserving record axis, to eliminate corresponding \(\mathsf{NULL}\) at the type level, must separately provide externalized carriers matching the three types of witnesses.
Among these: the worst-case and average budget lower bounds for collision direction are given by Theorem \ref{thm:address-defect-law}, so residue/inverse-limit axis cannot be omitted; reversibility of imaginary folding forces externalizing winding certificate axis (Theorem \ref{thm:dephys-dirichlet-winding-minimal}); covariance of endpoint continuous accumulation contraction is locked by uniqueness of Koenigs depth (Theorem \ref{thm:dephys-endpoint-koenigs-rigidity}); and renormalized finite part and universal pole of endpoint even channels provide an implementation-neutral closed-form sealing coordinate (Theorem \ref{thm:endpoint-even-finite-part-renormalization}).
\par
Under the "mutually non-coupling/canonical decomposition semantics" assumption, the above externalized components can be assembled in the product/direct sum sense as the product record axis \(r_0\), and from the universal property of products: any other extension-preserving record axis must be a quotient of this product; uniqueness of the quotient morphism comes from minimality and non-substitutability of each orthogonal component.
\end{proof}

\begin{corollary}[Unique continuous transverse direction: quotient uniqueness of radial register]\label{cor:xi-unique-continuous-transverse-register}
In the pure-phase regime of Definition \ref{def:xi-visible-read-null}, radial degrees of freedom deviating from the unitary slice are rejected by the type system and read out as \(\mathsf{NULL}\) (Theorem \ref{thm:xi-pw-no-continuous-hair}).
If one allows externalizing a \emph{continuous} mode axis to carry this radial degree of freedom for extension-preserving purposes, then this continuous axis must and can only fall, in the quotient sense, on the radial register
\[
\varrho(u):=\log|u|
\]
(calibration of \(\operatorname{proj}^{\mathrm e}\) in Remark \ref{rem:xi-soundness-vs-completeness}).
More precisely: any continuous record axis \(r:\CC^\times\to\RR\) that makes
\(
u\mapsto \bigl(\operatorname{proj}_{\mathsf{Phase}}(u),r(u)\bigr)
\)
injective must factor through \(|u|\), hence differs from \(\varrho\) only by a monotone reparametrization.
Therefore, there exists no second continuous transverse register parallel to the phase axis outside the unitary slice; any continuous extension not equivalent to \(\varrho\) must trigger \(\mathsf{NULL}\) in closure semantics.
\end{corollary}

\begin{proof}
If \(\bigl(\operatorname{proj}_{\mathsf{Phase}}(u),r(u)\bigr)\) is injective, then for any \(u,u'\) with the same phase \(\operatorname{proj}_{\mathsf{Phase}}(u)=\operatorname{proj}_{\mathsf{Phase}}(u')\), we must have \(r(u)\neq r(u')\) if and only if \(|u|\neq |u'|\);
therefore \(r\) gives an injection of \(r>0\) on each phase ray \(u=re^{i\theta}\), and by continuity this injection must be monotone.
Thus \(r\) must factor through the modulus map \(|\cdot|:\CC^\times\to\RR_{>0}\), and is equivalent (up to a monotone reparametrization) to the logarithmic coordinate \(\varrho=\log|\cdot|\) in the continuous injection sense.
The last sentence is given by "additional continuous coordinate registers cannot appear" in Theorem \ref{thm:xi-pw-no-continuous-hair}: if one refuses to externalize the above radial axis, then radial degrees of freedom are compiled into \(\mathsf{NULL}\) at the type level. \qed
\end{proof}

\begin{theorem}[Complete trichotomy of global \texorpdfstring{$\mathsf{NULL}$}{NULL} and interpretability (offline verifiable)]\label{thm:xi-null-complete-trichotomy-offline}
Under the visible domain constraints of Definition \ref{def:xi-visible-read-null}, the injectification principle of Proposition \ref{prop:xi-null-three-way}, and \(\mathbf{PW}\) closure type safety (Theorem \ref{thm:xi-offset-null-type-safety}),
for any target query/address \(q\) (which can be any protocol input object), if its global readout is \(\mathsf{NULL}\), then there must exist a \textbf{finite} failure witness \(\mathrm{Wit}(q)\) such that third parties can reproduce its correctness under offline, deterministic rules; and the type of this witness must fall into one of the following three categories:
\begin{enumerate}
  \item \textbf{Denotation/out-of-domain rejection (semantic or protocol null):} Address not given, or the interpretation of the object requires degrees of freedom beyond the visible domain \(X_{\mathrm{vis}}\) (typically: requires radial mode \(|u|\neq 1\) while visible domain is locked to \(\mathsf{Phase}\)), hence there is no address in closure and protocol reads out \(\mathsf{NULL}\);
  \item \textbf{Gluing knot (second-order compatibility obstruction):} Local certificate objects exist piecewise on the covering, but cannot be globalized due to second-order coherence constraint failure; its failure can be witnessed by a verifiable Čech \(2\)-cocycle obstruction class \([\alpha]\neq 0\) (Definition \ref{def:prefix-site-h2-gluing-obstruction}, Proposition \ref{prop:null-as-h2-obstruction});
  \item \textbf{Hard resource insufficiency (threshold-type budget violation):} To elevate many-to-one compression readout to single-valued verifiable readout, residue/precision budget fails to meet hard threshold, typically including:
  \begin{itemize}
    \item Modulus product budget of prime-register fails to cross quantile threshold (Proposition \ref{prop:xi-null-three-way});
    \item Endpoint layer resolution/precision budget insufficient to carry the quadratic compression gap of off-slice signal (Theorem \ref{thm:xi-offcritical-endpoint-resolution-lower-bound}; see also Theorem \ref{thm:xi-radial-information-projection-lower-bound}).
  \end{itemize}
\end{enumerate}
Moreover, the three types of witnesses are mutually orthogonal and non-substitutable within the system: increasing budget on one axis cannot substitute for fixing defects on another axis.
\end{theorem}

\begin{proof}
This trichotomy and non-substitutability is a restatement of the conclusion of Theorem \ref{thm:xi-null-orthogonal-trichotomy-resource-nonsubstitutable} under the calibration of this section.
The requirement "interpretation must be offline verifiable" comes from the certificate semantics of this paper: any acceptable assertion must be encoded as a finite object and reproduced by a deterministic verifier (see Definition \ref{def:xi-cert-refinement-equivalence} and deterministic readout in Theorem \ref{thm:xi-us-deterministic-tristate}), hence failure witnesses must also be output in the form of the same class of finite objects, otherwise they do not constitute interpretable assertions in protocol semantics.
\end{proof}

\begin{remark}[Closure defect witness when rejecting radial register: finite prime set and interval certificate]\label{rem:xi-null-closure-defect-finite-prime-witness}
Corollary \ref{cor:xi-unique-continuous-transverse-register} shows: if one wishes to carry off-slice information in continuous direction in extension-preserving manner, then the continuous part of mode axis can, in quotient sense, only be provided by radial register \(\varrho(u)=\log|u|\).
In any protocol that interprets finite-support data of residue/inverse-limit axis as logarithmic-scale ledger (e.g., records encoded by finite prime set \(S\) and integer vector \(v=(v_p)_{p\in S}\in\ZZ^{S}\)), the natural "scale closure" constraint can be written as
\[
\varrho_{\mathrm{vis}}(q)=\sum_{p\in S} v_p(q)\,\log p,
\]
where \(\varrho_{\mathrm{vis}}(q)\) is the visible readout of radial mode axis (if protocol refuses to externalize this axis, then pure-phase regime forces \(\varrho_{\mathrm{vis}}(q)=0\); see Theorem \ref{thm:xi-pw-no-continuous-hair}).
Therefore when rejecting radial register, closure defect is forced to be rewritten as a \textbf{finite} residue
\[
\varrho_{\mathrm{def}}(q):=\varrho_{\mathrm{vis}}(q)-\sum_{p\in S} v_p(q)\,\log p,
\]
whose nonzeroness can be reproduced offline by a dyadic interval certificate \([L_{\mathrm{def}},U_{\mathrm{def}}]\ni \varrho_{\mathrm{def}}(q)\) with \(0\notin [L_{\mathrm{def}},U_{\mathrm{def}}]\).
In the simplest case \(S=\{p\}\) and \(v_p(q)=1\), rejecting radial axis gives \(\varrho_{\mathrm{vis}}(q)=0\), hence
\(
\varrho_{\mathrm{def}}(q)=-\log p\neq 0
\),
which itself constitutes a witness of closure defect.
This "closure residue" is a standard auditable instance of category (3) (hard resource/closure defect) of Theorem \ref{thm:xi-null-complete-trichotomy-offline}.
\end{remark}

\paragraph{Precision coherence and implementation neutrality: inverse limit identity of certificate objects}
On the unitary slice, zero auditing can be compressed into event readout of a real-valued function on the phase axis: for example
\[
F(\theta):=\det\!\bigl(I-U_L(\theta)\bigr)
\]
(Proposition \ref{prop:xi-unitary-slice-spectral-crossing}).
For such readout, the "auditable" meaning adopted in this paper is: assertions must be given by finite data certificates and reproducible by deterministic offline rules (see Remark \ref{rem:xi-hard-zero-certificate} for explanation of univariate root interval certificate interface; also cf. Proposition \ref{prop:xi-rh-interval-criterion} and its Sturm/interval calibration implementation notes).

\begin{definition}[Certificate object, refinement, and verifiable equivalence]\label{def:xi-cert-refinement-equivalence}
Fix a verifiable real-valued function \(F:\RR\to\RR\) on the unitary slice, and fix a deterministic verifier \(\mathrm{Ver}\);
the verifier's input is a finite object \(C\) (encoded by finite integer/rational data), output is \(\mathsf{pass}/\mathsf{fail}\).
\par
For each dyadic depth \(b\in\NN\), let \(\mathsf{Cert}_b(F)\) denote the set of all certificate objects satisfying the following conditions:
\[
C=(I,\sigma),\qquad I=[\theta_{-},\theta_{+}]\subset\RR,\ \ \theta_{-},\theta_{+}\in 2^{-b}\ZZ,
\]
where \(I\) is an interval bracket, \(\sigma\) is auxiliary data (e.g., root interval certificate, single-root certificate, or other uniqueness/contraction certificates consistent with the protocol), and require
\(
\mathrm{Ver}(F;C)=\mathsf{pass}
\).
\par
Define \textbf{refinement} relation \(C'\preceq C\) (allowing \(b'\ge b\)) if and only if:
\begin{enumerate}
  \item \(I(C')\subseteq I(C)\);
  \item \(C'\) passes under the same verifier, i.e., \(\mathrm{Ver}(F;C')=\mathsf{pass}\).
\end{enumerate}
And define \textbf{verifiable equivalence} \(C\sim C'\) if and only if there exists a common refinement \(D\) such that \(D\preceq C\) and \(D\preceq C'\).
\end{definition}

\begin{theorem}[Inverse limit merger of precision coherence and implementation neutrality]\label{thm:xi-cert-inverse-limit-coherence}
Follow notation of Definition \ref{def:xi-cert-refinement-equivalence}.
Suppose there exists some \(\theta_\star\in\RR\) and a certificate chain \((C_b)_{b\ge b_0}\) satisfying:
\[
C_{b}\in\mathsf{Cert}_b(F),\qquad C_{b+1}\preceq C_{b},\qquad \bigl|I(C_b)\bigr|\xrightarrow[b\to\infty]{}0,
\]
and each \(C_b\)'s assertion isolates the same zero within \(I(C_b)\) in the sense of "at most one root" (e.g., Sturm root isolation or equivalent uniqueness certificate calibration, see Remark \ref{rem:xi-hard-zero-certificate}).
Then:
\begin{enumerate}
  \item \textbf{(Certificate identity)} Any valid certificate object denoting \(\theta_\star\) merges into a single equivalence class under \(\sim\);
  this equivalence class can be represented by the \emph{inverse limit} of a nested chain on the inverse system \((\mathsf{Cert}_b(F),\preceq)\), whose interval intersection is the singleton \(\bigcap_{b\ge b_0} I(C_b)=\{\theta_\star\}\).
  \item \textbf{(Implementation neutrality/no side information channel)}
  If verifier \(\mathrm{Ver}\) depends only on finite encoding of certificate objects (and operates according to fixed rules at integer/rational arithmetic layer),
  then for any two implementation environments, as long as the same certificate object \(C\) is input, verification output is strictly identical;
  and for fixed \(F\), addressing of \(\theta_\star\) depends only on its \(\sim\)-equivalence class, independent of non-certificate data such as floating-point trajectories or sampling paths that may arise during certificate generation.
\end{enumerate}
\end{theorem}

\begin{proof}
\begin{enumerate}
  \item
  By assumption, the "at most one root" certificate semantics ensures: if \(C,C'\) both denote the same zero \(\theta_\star\), then for sufficiently large \(b\) one can select a smaller certificate bracket refined to \(I(C)\cap I(C')\), hence there exists a common refinement \(D\);
  therefore any two valid certificates are equivalent under \(\sim\), merging into a single equivalence class.
  Nested interval length tending to \(0\) implies intersection is a single point, consistent with standard representation of inverse limit.
  \item
  If \(\mathrm{Ver}\) is deterministic computation on finite encoding, then its output is a function value of the input certificate object;
  hence output for the same \(C\) is invariant across different implementation environments.
  Meanwhile, by Definition \ref{def:xi-cert-refinement-equivalence}, equivalence classes depend only on the verifiable "can-be-commonly-refined" relation,
  and any generation details not entering certificate encoding are actively discarded at the semantic layer, hence do not affect addressing.
\end{enumerate}
\qed
\end{proof}

\begin{corollary}[Type negation: zero is not an address]\label{cor:xi-zero-not-address}
In the verifiable semantics of this paper (Definition \ref{def:xi-cert-refinement-equivalence}), a so-called "address object" is a \emph{finite} certificate object
\(
C=(I,\sigma)\in\mathsf{Cert}_b(F)
\)
(or its equivalence class under can-be-commonly-refined equivalence \(\sim\)), whose role is to specify an offline-reproducible assertion via finite encoding.
In contrast, Theorem \ref{thm:xi-cert-inverse-limit-coherence} shows: zero \(\theta_\star\) belongs to a \emph{point object} in the inverse limit sense of the inverse system \((\mathsf{Cert}_b(F),\preceq)\), whose denotation is given by interval intersection of an infinite refinement chain
\(
\bigcap_{b\ge b_0} I(C_b)=\{\theta_\star\}
\).
Therefore, "one zero is an address of another zero" does not hold at the type level: point objects cannot serve as finite certificate objects.
The statement that can hold and has verifiable meaning under the regime of this paper is: \emph{two zeros share the same visible prefix at given dyadic depth \(b\)} (hence are indistinguishable at that resolution), formalized by the dyadic prefix compatibility depth below.
\end{corollary}

\begin{definition}[Dyadic prefix compatibility depth and address ultrametric]\label{def:xi-dyadic-prefix-depth-ultrametric}
For any \(b\in\NN\) and \(\theta\in\RR\), define dyadic box
\[
J_b(\theta):=\Bigl[\frac{\lfloor 2^b\theta\rfloor}{2^b},\ \frac{\lfloor 2^b\theta\rfloor+1}{2^b}\Bigr),
\qquad |J_b(\theta)|=2^{-b}.
\]
For \(\theta,\theta'\in\RR\) define their dyadic prefix compatibility depth
\[
m(\theta,\theta'):=\sup\{\,b\in\NN:\ J_b(\theta)=J_b(\theta')\,\}\in\NN\cup\{\infty\},
\]
and agree \(2^{-\infty}=0\).
From this define \textbf{address ultrametric}
\[
d_{\mathrm{addr}}(\theta,\theta'):=2^{-m(\theta,\theta')}.
\]
Under this convention, \(d_{\mathrm{addr}}\) is an ultrametric, and its radius-\(2^{-b}\) closed ball is precisely the dyadic box: for any \(\theta\in\RR\),
\(
\overline{B}(\theta,2^{-b})=J_b(\theta)
\).
\end{definition}

\begin{remark}[Shared prefix and certificate refinement depth]\label{rem:xi-shared-prefix-cert-depth}
The \(m(\theta,\theta')\) of Definition \ref{def:xi-dyadic-prefix-depth-ultrametric} gives a fully verifiable "shared prefix" characterization: \(m(\theta,\theta')\ge b\) if and only if \(\theta,\theta'\) fall in the same dyadic box \(J_b\).
Therefore, for two zeros \(\theta_\star,\theta_\star'\), "sharing \(b\)-bit prefix" should be understood as: at the certificate type index \(b\) level (i.e., interval brackets with endpoints in \(2^{-b}\ZZ\)), they cannot yet be separated by any dyadic certificate of depth \(b\); but once refined to \(b>m(\theta_\star,\theta_\star')\), they must fall in different dyadic boxes, hence distinguishable at that depth.
\par
This ultrametric structure is isomorphic to the "ball equals cylinder set" calibration in the main text: the prefix ultrametric \(d_{\mathrm{pre}}\) induced by common prefix length on \(\Omega_\infty\) (paragraph \ref{par:spg-prefix-ultrametric}) and \(d_{\mathrm{addr}}\) here are merely dyadic realizations of the same "prefix refinement \(=\) radius contraction."
Therefore, in the address semantics of this paper, "resolution relation between zeros" should be understood as shared prefix depth (or equivalent ultrametric distance), rather than mistaking one point object as an address of another point object.
\end{remark}

\begin{remark}[Resolution scale and finite-state density barrier]\label{rem:xi-resolution-height-finite-state-barrier}
In the endpoint compactification/prefix semantics of this paper, "clarity" first manifests as a quantitative relation between height \(T\) and dyadic depth budget \(b\): in phase compactification coordinate \(u=\upsilon(T)\), the average distinguishable scale of adjacent zeros is \(|u_{n+1}-u_n|\asymp (T^{2}\log T)^{-1}\), hence the bit depth required to distinguish adjacent zeros satisfies
\[
b=2\log_{2}T+\log_{2}\log T+O(1)
\qquad(T\to\infty)
\]
(Appendix Proposition \ref{prop:app-adjacent-zero-dyadic-resolution}).
\par
On the other hand, any fixed finite-state/finite-syntax version of "zeros addressing each other" encounters density-order barrier: pole geometry of finite-dimensional determinant templates is locked to finitely many vertical lines and finite phase combs (Remark \ref{rem:finite-pole-lattice-barrier}), and in finite-band/finite-scale-window scale superposition regime, critical line zero count is at most \(N(T)=O(T)\) (Proposition \ref{prop:xi-scale-mellin-zero-density-upper}).
Therefore, to approach \(T\log T\) behavior at density order, the protocol must release exogenous degree-of-freedom budget growing with \(T\) in some direction (minimal mechanism is \(\Theta(\log T)\), see Remark \ref{rem:xi-scale-mellin-logT-budget}), and cannot be endogenously produced by self-referential addressing between zeros under fixed rules.
\end{remark}

\begin{theorem}[Verifier-program-layer repeatability and precision coherence: certificate existence partial function]\label{thm:xi-us-deterministic-tristate}
Under unitary slice audit calibration, fix:
\begin{enumerate}
  \item A visible domain constraint and checker (Definition \ref{def:xi-visible-read-null}, take \(X_{\mathrm{vis}}=\mathsf{Phase}\));
  \item A deterministic verifier \(\mathrm{Ver}\) and its certificate backend (Definition \ref{def:xi-cert-refinement-equivalence}), and agree certificate resolution/precision index is typed by dyadic depth \(b\in\NN\) (i.e., certificate endpoints fall in \(2^{-b}\ZZ\));
  \item A set of fixed implementation conventions (algorithm identifier, outward rounding and canonical input conventions), so that any "candidate certificate object" can be offline deterministically reproduced for \(\mathsf{pass}/\mathsf{fail}\) by third parties.
\end{enumerate}
Then there exists a deterministic three-valued partial function (according to typed address readout)
\[
\mathrm{Read}_{\mathrm{US}}:\ \mathrm{Addr}_{\mathrm{US}}\ \rightharpoonup\ \{\mathsf{NonZero},\mathsf{RootInterval},\mathsf{NULL}\},
\qquad
\mathrm{Addr}_{\mathrm{US}}:=\{(L,\theta,b):\ L>1,\ \theta\in\RR,\ b\in\NN\},
\]
with semantics: for each input address \((L,\theta,b)\), either output a finite interval certificate passing verification at depth \(b\) (e.g., one of two types "nonzero certificate/root interval certificate"; see Remark \ref{rem:xi-hard-zero-certificate} and Proposition \ref{prop:xi-rh-interval-criterion}), or output null \(\mathsf{NULL}\).
Moreover:
\begin{enumerate}
  \item \textbf{(Cross-implementation reproducibility)} Under fixed implementation conventions, \(\mathrm{Read}_{\mathrm{US}}\) is a deterministic partial function; same input address gives same output across different platforms/implementation environments.
  \item \textbf{(Object semantics of \(\mathsf{NULL}\))} \(\mathsf{NULL}\) does not express numerical instability, but structural absence of certificate objects: in site/(pre)sheaf calibration it is equivalent to "no protocol-allowed local section exists at this address/certificate fiber is empty" (Proposition \ref{prop:null-as-local-section-obstruction}); in concrete unitary certificate backend, it is equivalent to "input bit string incompatible with certificate backend hence no acceptable certificate" (Proposition \ref{prop:unitary-two-integer-reproducible}).
\end{enumerate}
\end{theorem}

\begin{proof}
Since \(\mathrm{Ver}\)'s operation on finite encoded objects is deterministic, any assertion "readout is non-\(\mathsf{NULL}\)" must be witnessed by some finite certificate object \(C\), and its verifiability is completely determined by truth value of \(\mathrm{Ver}(F;C)\) (Definition \ref{def:xi-cert-refinement-equivalence}).
Therefore under fixed implementation conventions, readout output depends only on input address and finite encoding of certificate object, hence cross-implementation consistent; this is precisely a direct corollary at readout layer of "implementation neutrality/no side information channel" in Theorem \ref{thm:xi-cert-inverse-limit-coherence}.
\par
On the other hand, Proposition \ref{prop:null-as-local-section-obstruction} intrinsicizes \(\mathsf{NULL}\) as emptiness of certificate fiber \(\mathscr F(a)\); while Proposition \ref{prop:unitary-two-integer-reproducible} gives "same bit string \(\Rightarrow\) unique certificate or unique null" replayability in a concrete, executable unitary certificate backend.
Together they ensure: source of \(\mathsf{NULL}\) is object-layer "non-denotable/non-constructible" (certificate does not exist), rather than floating-point trajectory or numerical noise. \qed
\end{proof}

\begin{corollary}[Formalization of no information leakage: typing precision index and explicit elevation]\label{cor:xi-precision-type-no-leakage}
Under the setting of Theorem \ref{thm:xi-us-deterministic-tristate}, dyadic depth \(b\) is elevated to certificate type index, and refinement relation \(C'\preceq C\) requires \emph{explicitly} entering deeper precision (Definition \ref{def:xi-cert-refinement-equivalence}).
Therefore, any operation attempting to combine/compare certificate objects in "mixed precision" without explicitly elevating to common depth and paying corresponding record cost is undefined in type semantics; its readout is equivalent by protocol to degrading to \(\mathsf{NULL}\).
In particular, this excludes the channel of smuggling additional side information into certificates through implicit mixed precision: all information that can affect verification conclusions must appear in the form of explicit precision index elevation and reproducible finite certificate objects.
\end{corollary}

\begin{proof}
Certificate objects are by definition organized hierarchically into \(\mathsf{Cert}_b(F)\), and any coherence of "same assertion" is only allowed to be expressed through refinement (\(b\) increasing) and can-be-commonly-refined equivalence relation \(\sim\) (Definition \ref{def:xi-cert-refinement-equivalence}; Theorem \ref{thm:xi-cert-inverse-limit-coherence}).
Therefore, if some process attempts to output a "mixed precision" candidate object without explicitly entering common depth, then it does not constitute a verifiable object in any fixed-\(b\) certificate category: either it cannot pass \(\mathrm{Ver}\), or it simply does not belong to protocol-allowed certificate fibers.
By Proposition \ref{prop:null-as-local-section-obstruction}, readout then manifests at object layer as \(\mathscr F(a)=\varnothing\), hence equivalent in protocol semantics to \(\mathsf{NULL}\).
Thus the conclusion holds. \qed
\end{proof}

\paragraph{Orthogonal slice nature of unitary slice: hidden axes can only appear as orthogonal invariants and externalized budgets}
The "orthogonality" of the unitary slice in this paper is not a Euclidean geometry proposition, but a structural direct sum characterized jointly by visible domain filtration, quotient decomposition, and endpoint payload decomposition.

\begin{theorem}[Orthogonal structure of critical line/unit circle slice: visible modulus component \(\oplus\) in-disk zero component \(\oplus\) endpoint payload component]\label{thm:xi-orthogonal-slice-structure}
Under unitary slice audit calibration of this paper (Definition \ref{def:xi-visible-read-null}, take \(X_{\mathrm{vis}}=\mathsf{Phase}\)) and \(\mathbf{PW}\) closure type checking:
\begin{enumerate}
  \item \textbf{(Protocol orthogonality: type rejection of radial mode)}
  Any object essentially requiring radial degree of freedom \(|u|\neq 1\) to express cannot be typed in pure-phase interface, its readout degrades by protocol to \(\mathsf{NULL}\) (Theorem \ref{thm:xi-offset-null-type-safety} and Theorem \ref{thm:xi-pw-no-continuous-hair}).
  \item \textbf{(Quotient-theoretic orthogonality: boundary modulus \(\perp\) in-disk zeros)}
  In Hardy/Smirnov condition calibration of disk completion \(F_\zeta\), boundary modulus determines outer factor, while in-disk zero data is carried by Blaschke inner factor;
  hence "deviation from critical line" information belongs to inner factor hidden component, structurally invisible to boundary modulus audit (\S\ref{subsec:unit-disk-inner-outer}).
  \item \textbf{(Endpoint payload orthogonality: modulo-\(4\) finite state \(\perp\) precision potential ledger)}
  In neighborhood of endpoint \(w=-1\), action of contraction semigroup decomposes into discrete endpoint state modulo \(4\) and additive ledger of precision potential \(y\), which do not couple (Proposition \ref{prop:endpoint-orthogonal-decomposition-z4-Rplus}).
\end{enumerate}
Therefore, "leaving \(\Re(s)=\tfrac12\)" if it is to remain verifiable in closed loop, can only appear along orthogonal invariants (defect ledger, endpoint payload) and externalized budget axes; if one refuses to externalize corresponding record axes, then \(\mathsf{NULL}\) must be triggered.
\end{theorem}
\begin{proof}
\begin{enumerate}
  \item
  Theorem \ref{thm:xi-offset-null-type-safety} and Theorem \ref{thm:xi-pw-no-continuous-hair} give protocol-layer type rejection: radial degrees of freedom requiring \(|u|\neq 1\) read out as \(\mathsf{NULL}\) in pure-phase interface.
  \item
  By inner-outer factor decomposition of \S\ref{subsec:unit-disk-inner-outer}: \(|F_\zeta(e^{i\theta})|\) determines only outer factor, in-disk zeros are encoded by inner factor, hence structurally separated in observable sense.
  \item
  Proposition \ref{prop:endpoint-orthogonal-decomposition-z4-Rplus} gives orthogonal decomposition "modulo-\(4\) finite state \(\oplus\) precision potential ledger" at endpoint.
\end{enumerate}
Merging three points yields the stated "orthogonal slice" structural conclusion. \qed
\end{proof}

\begin{corollary}[Orthogonal register necessity of inner factor data]\label{cor:xi-inner-factor-orthogonal-register}
Under the setting of Theorem \ref{thm:xi-orthogonal-slice-structure}, further assume Hardy/Smirnov condition calibration of Appendix \ref{subsec:unit-disk-inner-outer} holds, so disk completion object has standard factorization
\(
F_\zeta(w)=B(w)\,S(w)\,O(w)
\).
Then:
\begin{enumerate}
  \item \textbf{(Off-critical-line information equals inner factor data)}
  Any off-critical-line zero \(\rho=\beta+\mathrm{i}\gamma\) (\(0<\beta<\tfrac12\)) has disk image \(w_\rho\in\mathbb{D}\) and its offset coordinate
  \(
  \delta:=\tfrac12-\beta=\mathcal{B}_{-1}(w_\rho)
  \)
  belongs to zero data of Blaschke inner factor \(B\); hence it is not determined by outer factor in unitary slice's \emph{boundary modulus} audit calibration.
  \item \textbf{(No orthogonal carrier then not rigorizeable)}
  If verification protocol maintains \(X_{\mathrm{vis}}=\mathsf{Phase}\) and refuses to explicitly externalize radial record on mode axis \(a_{\mathrm{mode}}\), then any strategy attempting to rigorize the above inner factor data as constant-layer verifiable objects of unitary slice will inevitably trigger protocol-layer type rejection: its readout either stably degrades to \(\mathsf{NULL}\), or must introduce additional orthogonal record axes to carry this inner factor data (consistent with falsifiable interface of Proposition \ref{prop:xi-offcritical-falsifiable-restatement}).
\end{enumerate}
\end{corollary}

\begin{proof}
Item (1) is direct correspondence of "in-disk zeros \(\Leftrightarrow\) Blaschke data" in Appendix \ref{subsec:unit-disk-inner-outer}, consistent with item (2) of Theorem \ref{thm:xi-orthogonal-slice-structure}.
Item (2) is derived from item (1) of Theorem \ref{thm:xi-orthogonal-slice-structure} (radial mode cannot be typed in pure-phase interface) and visible domain filtration "address precedes value" (Definition \ref{def:xi-visible-read-null}); when verifiable completeness is required, the unique non-degenerate path is to externalize corresponding orthogonal record axes. \qed
\end{proof}

\paragraph{\texorpdfstring{$\Xi_\zeta$}{Xi-zeta}-type zero statistics logarithmic degree-of-freedom lower bound: linear density barrier and Hankel refuter}
This paragraph compresses the comb rigidity of "finite-dimensional/finite-band" and "\(\log T\) degree-of-freedom" threshold into a directly refutable statistical lower bound:
if the goal is to approximate the classical \(\zeta\) critical line zero count main order \(N_\zeta(T)\asymp T\log T\), then any implementation calibration that remains long-term at fixed finite bandwidth, fixed finite scale window, or uniformly bounded Hankel dimension will structurally fail at density order.

\begin{theorem}[\texorpdfstring{$\Xi_\zeta$}{Xi-zeta}-type zero statistics logarithmic degree-of-freedom lower bound (verifiable refuter)]\label{thm:xi-zero-statistics-log-budget-lower-bound}
Within the completion/\(\Xi\) audit framework of this paper, the following three structural necessary conditions hold (all are auditable/computable refutation interfaces):
\begin{enumerate}
  \item \textbf{(Linear density barrier)}
  If scale superposition completion \(\mathcal X_{\Delta,w}\) (Definition \ref{def:xi-scale-mellin-superposition}) kernel \(\Delta(z,u)\) has finite bandwidth on unitary slice, and weight function \(w\) has finite scale window (compact support) on logarithmic scale, then critical line section
  \(t\mapsto \mathcal X_{\Delta,w}(\tfrac12+it)\)
  real zero count satisfies linear upper bound \(N(T)=O(T)\) (Proposition \ref{prop:xi-scale-mellin-zero-density-upper}).
  Hence this "finite-band + finite scale window" regime is incompatible with \(T\log T\) at density order.

  \item \textbf{(\(\log T\) degree-of-freedom threshold)}
  If goal is to approximate \(N_\zeta(T)\asymp T\log T\), then must release \(\Theta(\log T)\)-level degree-of-freedom budget in some orthogonal direction:
  either effective scale window width \(X(T)\) grows with \(T\) to \(X(T)\asymp \log T\), or critical circle mode count \(B(T)\) grows with \(T\) to \(B(T)\asymp (\log T)/(\log\lambda)\)
  (Remark \ref{rem:xi-scale-mellin-logT-budget}; Corollary \ref{cor:critical-circle-mode-growth}).
  In particular, in fixed finite-dimensional kernel single-scale Dirichletization calibration, zero/pole count is at most \(O(T)\), hence must introduce mode/state count growing with \(T\) to cross this barrier (Corollary \ref{cor:single-scale-dirichlet-lattice} and Remark \ref{rem:single-scale-dirichlet-obstruction}).

  \item \textbf{(Hankel complexity refuter)}
  For any auditable backend, if its visible output sequence (e.g., trace sequence or finite window segment thereof) maintains consistently bounded Hankel rank in the limit of elevating height/resolution,
  then its minimal linear realization dimension remains bounded in protocol sense (Theorem \ref{thm:hankel-rank-minimal}).
  Since this dimension is a hard lower bound on "required visible degrees of freedom" (Remark \ref{rem:xi-acc-rank}), consistently bounded Hankel rank is incompatible with the \(\log T\) threshold of item (2);
  hence "Hankel rank does not grow" can serve as an independent, computable structural refuter: it refutes any claim of "isomorphic to \(\Xi_\zeta\) at density order."
\end{enumerate}
\end{theorem}

\begin{proof}
Item (1) is Proposition \ref{prop:xi-scale-mellin-zero-density-upper}.
Item (2) is a threshold corollary of Proposition \ref{prop:xi-scale-mellin-zero-density-upper}: if maintaining finite scale window and finite bandwidth, then zero count is at most linear; hence to approximate \(T\log T\) must break one of "finite-band/finite scale window," and under minimal mechanism introduce \(\Theta(\log T)\) degrees of freedom (Remark \ref{rem:xi-scale-mellin-logT-budget}).
In finite-dimensional critical circle comb framework, calibration of critical line density and mode count \(B\) is given by Corollary \ref{cor:critical-circle-mode-growth}.
Item (3) is directly derived from Hankel rank theorem (Theorem \ref{thm:hankel-rank-minimal}) and ACC-Rank stability calibration (Remark \ref{rem:xi-acc-rank}): Hankel rank is certificated lower bound on minimal realization dimension; if bounded in limit, then "required visible degrees of freedom" remains bounded, contradicting item (2)'s requirement of \(\log T\)-level degree-of-freedom release. \qed
\end{proof}

\paragraph{Visibility limit of symmetry in disk interface: quadratic radial compression law of off-slice offset}
In appendix disk completion interface, "leaving critical slice \(\Re(s)=\tfrac12\)" is rigorously transported as point-like Blaschke data inside unit disk \(\mathbb{D}\).
Under fixed visible interface (involution-fixed slice and parameter-free compactification gate), this off-slice information can only enter through second-order radial channel in endpoint neighborhood (Corollary \ref{cor:xi-null-second-order-radial-channel}).

\subparagraph{Structural source of endpoint compression and precise formulation of "scale transport" problem}
Following horizon notation of \S\ref{subsubsec:xi-ramanujan-horizon-ledger-framework} and Appendix \ref{subsubsec:app-horizon-weyl-carath-toeplitz}, let
\[
E_{\zeta}(z):=\Xi_{\zeta}\!\left(\tfrac12+\mathrm{i}z\right),\qquad
m_{\zeta}(z):=-\frac{E_{\zeta}'(z)}{E_{\zeta}(z)}\qquad(z\in\mathbb{H}),
\]
and let \(w=\mathcal{C}(z)\) be Cayley compactification, \(s(w)=\tfrac12+\mathrm{i}\,\mathcal{C}^{-1}(w)\) be disk section, so disk horizon logarithmic derivative
\[
\mathscr{C}_{\zeta}(w):=-\frac{\Xi_{\zeta}'(s(w))}{\Xi_{\zeta}(s(w))}\qquad(w\in\mathbb{D})
\]
satisfies strict identity \(\mathscr{C}_{\zeta}(w)=-\mathrm{i}\,m_{\zeta}(\mathcal{C}^{-1}(w))\) (equation \eqref{eq:app-horizon-m-c-relation}).
In the dephysicalized semantics of this paper, \(m_{\zeta}\) (equivalently \(\mathscr{C}_{\zeta}\)) can be viewed as horizon impedance/Weyl function, and its meromorphic poles in \(\mathbb{H}\) (equivalently \(\mathbb{D}\)) defined as \textbf{horizon resonance poles}.

\begin{proposition}[Positivity criterion of extreme resonance and finite-dimensional failure witness]\label{prop:xi-extreme-resonance-positivity}
Defining resonance as meromorphic poles of \(m_{\zeta}\) (or equivalently \(\mathscr{C}_{\zeta}\)) in \(\mathbb{H}\) (or \(\mathbb{D}\)):
\begin{enumerate}
  \item \textbf{(Off-slice \texorpdfstring{$\Longleftrightarrow$}{<->} upper-half-plane/disk interior pole)}
  If \(\Xi_{\zeta}\) has off-critical-line zero \(\rho=\beta+\mathrm{i}\gamma\) (\(\beta<\tfrac12\)), then
  \(z=\gamma+\mathrm{i}(\tfrac12-\beta)\in\mathbb{H}\) satisfies \(E_{\zeta}(z)=0\), hence \(m_{\zeta}\) has pole in \(\mathbb{H}\); equivalently, \(\mathscr{C}_{\zeta}\) has pole at disk image \(w=\mathcal{C}(z)\in\mathbb{D}\) (unstable resonance).
  \item \textbf{(RH \texorpdfstring{$\Longleftrightarrow$}{<->} positivity \texorpdfstring{$\Longleftrightarrow$}{<->} poles pushed to boundary)}
  RH holds if and only if \(m_{\zeta}\) is Herglotz function (Theorem \ref{thm:app-horizon-weyl-herglotz}), also if and only if \(\mathscr{C}_{\zeta}\) belongs to Carathéodory class (Corollary \ref{cor:app-horizon-carath-rh}).
  Hence under RH, upper-half-plane/disk interior resonance poles are structurally excluded: allowed singularities can only be carried as boundary spectral measure (Herglotz/Carathéodory representation; see Remark \ref{rem:app-horizon-multipole-positivity-hair}).
  \item \textbf{(Toeplitz--PSD certificatization)}
  Carathéodory positivity further equivalent to Toeplitz--PSD hierarchy \(\{\mathbf{T}_N\succeq 0\}_{N\ge 0}\) (Theorem \ref{thm:app-horizon-toeplitz-lmi}).
  Therefore once in-disk pole appears (equivalent to off-slice zero), there must exist some finite level \(N\) such that \(\mathbf{T}_N\) has negative eigenvalue, giving offline-reproducible finite-dimensional failure witness.
  Furthermore, coefficient data of Toeplitz hierarchy has complementary auditable sources: in absolutely convergent patch \(|w|<1/3\), \(\mathscr{C}_{\zeta}\) can be term-by-term audited by Euler product/Dirichlet series (Theorem \ref{thm:app-horizon-euler-patch}); while in local calibration, \(\mathscr{C}_{\zeta}\) coefficients can also be explicitly generated by finite-order logarithmic derivative locally sealed at \(s=-\tfrac12\) (Proposition \ref{prop:app-horizon-coeff-local}).
\end{enumerate}
\end{proposition}

\begin{proof}
Three assertions directly derived from Theorem \ref{thm:app-horizon-weyl-herglotz}, Corollary \ref{cor:app-horizon-carath-rh}, Theorem \ref{thm:app-horizon-toeplitz-lmi} and identity \eqref{eq:app-horizon-m-c-relation}. \qed
\end{proof}

\subparagraph{Schurization of Carathéodory impedance and Schur--OPUC certificate chain}
Proposition \ref{prop:xi-extreme-resonance-positivity} has strictly aligned "appearance of resonance poles in stability domain" with failure of Carathéodory positivity and Toeplitz--PSD hierarchy.
On top of this positivity interface, Cayley--Schur transformation further rewrites positive-real-part property of horizon impedance as Schur class, thereby elevating "resonance=pole" dictionary to a fully discrete reflection coefficient certificate chain, and giving OPUC closed-form readout of endpoint threshold carrying strength.
\par
Define Cayley--Schur transformation
\begin{equation}\label{eq:xi-horizon-cayley-schur}
\boxed{\;
S_{\zeta}(w):=\frac{\mathscr{C}_{\zeta}(w)-1}{\mathscr{C}_{\zeta}(w)+1},
\qquad
\mathscr{C}_{\zeta}(w)=\frac{1+S_{\zeta}(w)}{1-S_{\zeta}(w)}.\;}
\end{equation}

\begin{theorem}[Carathéodory \texorpdfstring{$\Longleftrightarrow$}{<->} Schur and stability domain pole exclusion]\label{thm:xi-horizon-carath-schur-equivalence}
Let \(\mathscr{C}_{\zeta}\) be analytic in \(\mathbb{D}\) and satisfy \(\Re \mathscr{C}_{\zeta}(w)\ge 0\).
Then \(S_{\zeta}\) is analytic in \(\mathbb{D}\) and satisfies \(|S_{\zeta}(w)|\le 1\), i.e., \(S_{\zeta}\) belongs to Schur class; conversely also holds.
Therefore when \(\mathscr{C}_{\zeta}\) is Carathéodory, stability domain \(\mathbb{D}\) interior resonance poles are structurally excluded, allowed singularities can only be carried as boundary spectral measure on \(|w|=1\) (Herglotz/Carathéodory representation).
\end{theorem}

\begin{definition}[Schur algorithm and reflection coefficients]\label{def:xi-horizon-schur-parameters}
Apply Schur algorithm to Schur function \(S_{\zeta}\), define
\[
S_0:=S_{\zeta},\qquad \alpha_n:=S_n(0),\qquad
S_{n+1}(w):=\frac{1}{w}\,\frac{S_n(w)-\alpha_n}{1-\overline{\alpha_n}S_n(w)}\qquad(n\ge 0).
\]
Sequence \((\alpha_n)_{n\ge 0}\) called Schur/Verblunsky (reflection) coefficients.
\end{definition}

\begin{theorem}[Finite refuter of reflection coefficient violation]\label{thm:xi-horizon-reflection-finite-witness}
If \(\mathscr{C}_{\zeta}\) is Carathéodory (equivalent \(S_{\zeta}\) is Schur), then for all \(n\ge 0\) have \(|\alpha_n|<1\).
Conversely, if there exists some \(N\) such that \(|\alpha_N|\ge 1\), then \(S_{\zeta}\) not Schur, hence \(\mathscr{C}_{\zeta}\) not Carathéodory;
therefore stability domain must have resonance pole, equivalently triggering some finite-level Toeplitz--PSD failure witness (item (3) of Proposition \ref{prop:xi-extreme-resonance-positivity}; see \cite{Simon2005OPUC1}).
\end{theorem}

\begin{theorem}[Conservative realization and endpoint threshold carrying]\label{thm:xi-horizon-conservative-realization}
If \(\mathscr{C}_{\zeta}\) is Carathéodory, then \(S_{\zeta}\) has conservative (unitary) scattering realization: there exists unitary colligation \(\mathcal U\) such that \(S_{\zeta}\) is its characteristic function.
Equivalently, there exists (finite/infinite) CMV canonical form realization generated by \((\alpha_n)\), and in this conservative phase stability domain has no scattering poles; all "extreme resonances" can only degenerate to threshold carrying on boundary \(|w|=1\) (see \cite{CanteroMoralVelazquez2003CMV,Simon2005OPUC1}, isomorphic to finite-window CMV calibration of Proposition \ref{prop:xi-hilbert-polya-cmv}).
\end{theorem}

\begin{theorem}[Christoffel formula for endpoint atomic strength]\label{thm:xi-horizon-endpoint-atom-christoffel}
Let \(\mathscr{C}_{\zeta}\) be Carathéodory, and let \(\mu_{\zeta}\) be its representing measure (Theorem \ref{thm:app-horizon-weyl-herglotz}, Corollary \ref{cor:app-horizon-carath-rh}).
Denote \(\{\varphi_n\}_{n\ge 0}\) unit circle orthonormal polynomials (OPUC) for \(\mu_{\zeta}\), and define Christoffel function
\[
\lambda_N(\xi):=\left(\sum_{n=0}^{N}|\varphi_n(\xi)|^2\right)^{-1},\qquad \xi\in\TT.
\]
Then for any \(\xi\in\TT\) holds
\[
\mu_{\zeta}(\{\xi\})
=\left(\sum_{n=0}^{\infty}|\varphi_n(\xi)|^2\right)^{-1}
=\lim_{N\to\infty}\lambda_N(\xi),
\]
and \(\mu_{\zeta}(\{\xi\})=0\) if and only if \(\lambda_N(\xi)\to 0\) (see \cite{Simon2005OPUC1}).
Take \(\xi=-1\) to get endpoint point-mass threshold strength
\[
\mathfrak m_{-1}:=\mu_{\zeta}(\{-1\})=\lim_{N\to\infty}\lambda_N(-1).
\]
\end{theorem}

\begin{corollary}[Two-stage criterion: reflection coefficient non-violation \(\times\) endpoint atom controllability]\label{cor:xi-horizon-extreme-resonance-two-stage}
Under Carathéodory/Toeplitz--PSD calibration, certificatization of extreme resonance can be decomposed into two complementary constraint groups:
\begin{enumerate}
  \item \textbf{Stability domain no resonance:} Schur parameters satisfy \(|\alpha_n|<1\); any violation \(|\alpha_N|\ge 1\) gives finite refuter (Theorem \ref{thm:xi-horizon-reflection-finite-witness}).
  \item \textbf{Boundary threshold carrying:} Atomic strength of endpoint \(-1\) precisely given by Christoffel limit (Theorem \ref{thm:xi-horizon-endpoint-atom-christoffel}), thereby providing recursively auditable point-wise quantization coordinate for endpoint carrying structure.
\end{enumerate}
\end{corollary}

Under the above positivity—stability semantics, Poisson/Busemann coordinate \(\mathcal{B}_{-1}\) of endpoint \(-1\) gives intrinsic readout scale of off-slice offset in disk interface; the following theorem gives complete closed form and high-height asymptotics of disk image \(w_\rho\), and attributes "visibility resolution" precisely to quadratic compression induced by precision potential \(\mathcal{P}\).
This compression simultaneously poses an intrinsic scale transport problem: if maintaining interface radius and certificate semantics unchanged, then distinguishable scale of high-height off-slice data in endpoint boundary layer contracts quadratically with \(|\gamma|\); to transport this information to fixed radius to form height-independent Jensen/PSD certificate, must introduce externalized scale gates, rigorous implementation see \S\ref{subsubsec:xi-adams-norm-scale-transport}.

\begin{theorem}[Quadratic radial compression law of off-critical-line offset and precision potential attribution]\label{thm:xi-offcritical-quadratic-radial-compression}
Assume there exists off-critical-line zero \(\rho=\beta+\mathrm{i}\gamma\) satisfying \(0<\beta<\tfrac12\).
Let
\(
\delta:=\tfrac12-\beta>0
\)
and let \(w_\rho=\mathcal{C}(\gamma+\mathrm{i}\delta)\in\mathbb{D}\) be its disk image.
Then:
\begin{enumerate}
  \item \textbf{(Closed-form identity/horocycle realization)}
  \[
  \delta=\mathcal{B}_{-1}(w_\rho)=\frac{1-|w_\rho|^2}{|1+w_\rho|^2},
  \qquad
  |w_\rho|^2=\frac{\gamma^2+(1-\delta)^2}{\gamma^2+(1+\delta)^2},
  \qquad
  1-|w_\rho|^2=\frac{4\delta}{\gamma^2+(1+\delta)^2}=\delta\,|1+w_\rho|^2
  \]
  (Corollary \ref{cor:app-offcritical-horocycle-busemann}).
  Therefore when \(|\gamma|\to\infty\) with \(\delta\) bounded,
  \[
  1-|w_\rho|^{2}=\frac{4\delta}{\gamma^{2}}\bigl(1+o(1)\bigr),
  \qquad
  1-|w_\rho|=\frac{2\delta}{\gamma^{2}}\bigl(1+o(1)\bigr),
  \]
  so at fixed offset \(\delta\), off-slice information is compressed into distinguishable scale \(\Theta(\gamma^{-2})\) of disk boundary layer, whose compression strength intensifies quadratically with height.
  This law is not empirical geometry in "symmetry axis" sense, but structural consequence of "involution-fixed slice \(|u|=1\) + compactification gate" in visibility calibration.
  \item \textbf{(Endpoint second-order radial channel)}
  Denote endpoint depth \(\delta_\gamma:=\tfrac12-\pi^{-1}\arctan|\gamma|\downarrow 0\), then
  \[
  1-|w_\rho|^2
  =4\pi^2\delta\,\delta_\gamma^2\bigl(1+o(1)\bigr)\qquad(\delta_\gamma\downarrow 0)
  \]
  (Corollary \ref{cor:app-offcritical-radius-second-order-radial-channel}).
  Hence in natural scale \(|1+w|\ll 1\) of endpoint neighborhood, \(\delta\) enters as \emph{quadratic} term in chordal distance \(|1+w|^2\).
  \item \textbf{(Precision potential attribution)}
  Let precision potential \(\mathcal{P}(x)=\log\pi+\log(1+x^2)\) (Definition \ref{def:unit-circle-phase-precision-potential}), then
  \[
  1-|w_\rho|^2\sim 4\pi\,\delta\,e^{-\mathcal{P}(\gamma)}\qquad(|\gamma|\to\infty,\ \delta\ \text{bounded})
  \]
  (Remark \ref{rem:app-offcritical-radius-precision-potential}).
  This attributes "disk visibility compression of off-critical-line offset" precisely to exponential factor \(e^{-\mathcal{P}(\gamma)}\) of precision potential.
\end{enumerate}
\end{theorem}

\begin{proof}
Three assertions are direct restatements of Corollary \ref{cor:app-offcritical-horocycle-busemann}, Corollary \ref{cor:app-offcritical-radius-second-order-radial-channel} and Remark \ref{rem:app-offcritical-radius-precision-potential}. \qed
\end{proof}

\begin{corollary}[Precision potential decomposition and "offset load": logarithmic additive structure of endpoint layer depth]\label{cor:xi-offcritical-endpoint-depth-log-decomposition}
Under setting of Theorem \ref{thm:xi-offcritical-quadratic-radial-compression}, denote endpoint layer margin
\(
h(\rho):=1-|w_\rho|^2
\).
Then when \(|\gamma|\to\infty\) with \(\delta\) bounded, have product-type asymptotics
\begin{equation}\label{eq:xi-offcritical-endpoint-product-law}
\boxed{\;
\frac{h(\rho)}{\delta}\,e^{\mathcal{P}(\gamma)}\to 4\pi\;}
\end{equation}
and equivalent logarithmic depth decomposition
\begin{equation}\label{eq:xi-offcritical-endpoint-depth-load-decomposition}
\boxed{\;
-\log h(\rho)=\mathcal{P}(\gamma)+\log\frac{1}{\delta}-\log(4\pi)+o(1)\qquad(|\gamma|\to\infty)\; }.
\end{equation}
\end{corollary}

\begin{proof}
From item (3) asymptotics of Theorem \ref{thm:xi-offcritical-quadratic-radial-compression}
\(
h(\rho)\sim 4\pi\,\delta\,e^{-\mathcal{P}(\gamma)}
\)
immediately get \eqref{eq:xi-offcritical-endpoint-product-law}.
Taking negative logarithm on both sides and rearranging yields \eqref{eq:xi-offcritical-endpoint-depth-load-decomposition}. \qed
\end{proof}

\begin{remark}[Comoving horizon localization: linear development and regime cost of address prefix]\label{rem:xi-comoving-horizon-localization-address-null}
The endpoint depth decomposition of equations \eqref{eq:xi-offcritical-endpoint-product-law}--\eqref{eq:xi-offcritical-endpoint-depth-load-decomposition} is obtained under fixed Cayley chart \(\mathcal{C}\) (i.e., fixed-basepoint radial audit): when \(|\gamma|\to\infty\) with \(\delta\) bounded, \(\delta\) is structurally compressed into second-order radial channel of endpoint neighborhood through factor \(e^{-\mathcal{P}(\gamma)}\asymp \gamma^{-2}\) (Theorem \ref{thm:xi-offcritical-quadratic-radial-compression}).
\par
However this "second-order compression" is not an intrinsic obstruction of horocycle structure.
Horizontal translation \(t\mapsto t-x_0\) (\(x_0\in\RR\)) in upper half-plane preserves \(\Im(t)\) invariant, hence manifests on disk endpoint Busemann coordinate \(\delta=\mathcal{B}_{-1}(w)\) (Corollary \ref{cor:app-offcritical-horocycle-busemann}) as disk automorphism reselection along same horocycle.
Appendix Theorem \ref{thm:app-comoving-horizon-localization} gives its closed-form effect: taking comoving parameter \(x_0=\gamma\),
\[
1-|w_{\rho,\gamma}|^2=\frac{4\delta}{(1+\delta)^2}\sim 4\delta\qquad(\delta\downarrow 0),
\]
so in comoving localization chart off-critical-line offset manifests at linear scale, no longer producing second-order compression with \(|\gamma|\).
\par
This localization geometrically dissolves second-order compression, but is not cost-free in verifiable regime: it is compatible with the "address precedes value" regime of this paper, and translation parameter \(x_0\) must be explicitly given as addressable prefix (\S\ref{sec:recursive-addressing}); when address not given, corresponding evaluation function is undefined at object layer and reads out as \(\mathsf{NULL}\), and recursive addressing only reselects and organizes events within existing visible domain without expanding visible \(\sigma\)-algebra (Corollary \ref{cor:recursive-addressing-visible-sigma-nonexpansion}).
Therefore, subsequent lower bounds on endpoint resolution and bit complexity (e.g., Theorem \ref{thm:xi-offcritical-endpoint-resolution-lower-bound}) still characterize unavoidable budget under fixed visible interface: if not externalizing this comoving prefix, then off-line offset in high-height intervals must enter through second-order radial channel and trigger corresponding capacity threshold.
\end{remark}

\begin{remark}[Additive load of endpoint depth]\label{rem:xi-offcritical-endpoint-depth-additive-load}
Equation \eqref{eq:xi-offcritical-endpoint-depth-load-decomposition} shows: under endpoint compression scale, endpoint layer depth \(-\log(1-|w_\rho|^2)\) decomposes asymptotically into two additive loads: height precision load \(\mathcal{P}(\gamma)\) induced by parameter-free compactification Jacobian, and offset load \(\log(1/\delta)\) contributed by off-critical-line deviation.
In verifiable semantics of this paper, this decomposition is compatible with "orthogonal slice" conclusion: radial information deviating from slice does not enter as freely adjustable additional continuous coordinate, but appears in form of auditable ledger load at endpoint layer.
\end{remark}

\begin{theorem}[Endpoint resolution—bit complexity barrier: irreducible budget lower bound derived from quadratic compression (explicit worst case)]\label{thm:xi-offcritical-endpoint-resolution-lower-bound}
Fix \(\delta_0\in(0,\tfrac12)\) and \(T>0\).
If there exists off-critical-line zero \(\rho=\beta+\mathrm{i}\gamma\) satisfying
\[
|\gamma|\le T,\qquad \beta\le \tfrac12-\delta_0,
\]
and still denote its disk image as \(w_\rho=\mathcal{C}(\gamma+\mathrm{i}(\tfrac12-\beta))\), then endpoint layer margin satisfies explicit lower bound
\begin{equation}\label{eq:xi-offcritical-endpoint-margin-min}
\boxed{\;
1-|w_\rho|^2\ \ge\ \frac{4\delta_0}{T^2+(1+\delta_0)^2}\; }.
\end{equation}
Therefore, any verification strategy attempting to exclude all off-line zeros in above height window solely by endpoint layer resolution threshold control, if its minimal distinguishable scale for margin \(1-|w|^2\) is \(2^{-p}\), must satisfy
\begin{equation}\label{eq:xi-offcritical-endpoint-bit-lower-bound}
\boxed{\;
p\ \ge\ \log_2\!\Bigl(\frac{T^2+(1+\delta_0)^2}{4\delta_0}\Bigr)\; }.
\end{equation}
And when \(T\to\infty\) have asymptotic lower bound
\begin{equation}\label{eq:xi-offcritical-endpoint-bit-asymptotic}
\boxed{\;
p\ \ge\ 2\log_2 T+\log_2\!\frac{1}{\delta_0}+O(1)\; }.
\end{equation}
Using precision potential \(\mathcal{P}(x)=\log\pi+\log(1+x^2)\), this asymptotics can also be written as
\begin{equation}\label{eq:xi-offcritical-endpoint-bit-precision-potential-form}
\boxed{\;
p\ \ge\ \frac{\mathcal{P}(T)+\log(1/\delta_0)-\log(4\pi)}{\log 2}+o(1)\qquad(T\to\infty)\; }.
\end{equation}
\end{theorem}

\begin{proof}
By closed-form identity of item (1) of Theorem \ref{thm:xi-offcritical-quadratic-radial-compression},
\[
1-|w_\rho|^2=\frac{4\delta}{\gamma^2+(1+\delta)^2},
\qquad \delta=\tfrac12-\beta\in[\delta_0,\tfrac12).
\]
On \(\delta\in(0,1)\), function \(\delta\mapsto \frac{4\delta}{\gamma^2+(1+\delta)^2}\) is monotonically increasing; and for \(|\gamma|\le T\) its right-hand side decreases as \(|\gamma|\) increases.
Therefore under constraints \(\delta\ge\delta_0,|\gamma|\le T\), minimum of margin is attained at \(\delta=\delta_0\) and \(|\gamma|=T\), hence \eqref{eq:xi-offcritical-endpoint-margin-min}.
\par
If verification mechanism's minimal distinguishable scale for endpoint margin is \(2^{-p}\), then to achieve endpoint-layer-distinguishable exclusion of all possible off-line zeros in this height window, must satisfy \(2^{-p}<\inf(1-|w_\rho|^2)\), combined with \eqref{eq:xi-offcritical-endpoint-margin-min} gives \eqref{eq:xi-offcritical-endpoint-bit-lower-bound}.
\eqref{eq:xi-offcritical-endpoint-bit-asymptotic} obtained from \(\log_2(T^2+(1+\delta_0)^2)=2\log_2 T+O(1)\).
Finally, from \(\mathcal{P}(T)=\log\pi+\log(1+T^2)\) can write main term in precision potential form, yielding \eqref{eq:xi-offcritical-endpoint-bit-precision-potential-form}. \qed
\end{proof}

\begin{corollary}[Bit budget gain law of comoving prefix: second-order compression \texorpdfstring{$\Rightarrow$}{=>} constant budget]\label{cor:xi-comoving-prefix-bit-budget-gain}
Follow setting of Theorem \ref{thm:xi-offcritical-endpoint-resolution-lower-bound} and fix \(\delta_0\in(0,\tfrac12)\).
If allowing object-layer externalization of addressable prefix \(x_0\) and taking comoving localization \(x_0=\gamma\) (Appendix Theorem \ref{thm:app-comoving-horizon-localization}; see Remark \ref{rem:xi-comoving-horizon-localization-address-null} for regime cost explanation), then off-line offset detectable margin in comoving chart satisfies
\[
h_{\gamma}(\rho)=1-|w_{\rho,\gamma}|^2=\frac{4\delta}{(1+\delta)^2}\ \ge\ \frac{4\delta_0}{(1+\delta_0)^2}\qquad(\delta\ge\delta_0).
\]
Therefore, if verification mechanism uses threshold \(2^{-p}\) as "detectable margin" (i.e., require \(h_{\gamma}(\rho)\ge 2^{-p}\) for all \(\delta\ge\delta_0\)), must satisfy height-independent constant lower bound
\begin{equation}\label{eq:xi-comoving-bit-threshold}
\boxed{\;
p\ \ge\ \log_2\!\Bigl(\frac{(1+\delta_0)^2}{4\delta_0}\Bigr)\; }.
\end{equation}
Comparing with \eqref{eq:xi-offcritical-endpoint-bit-asymptotic}: in height window \(|\gamma|\le T\), main-order bit budget "recovered" by comoving prefix is \(2\log_2T+O(1)\).
\end{corollary}

\begin{proof}
By Appendix Theorem \ref{thm:app-comoving-horizon-localization}, when taking \(x_0=\gamma\)
\(
h_{\gamma}(\rho)=\frac{4\delta}{(1+\delta)^2}
\).
On \(\delta\in(0,1)\) this function is monotonically increasing, so for all \(\delta\ge\delta_0\) have
\(
h_{\gamma}(\rho)\ge \frac{4\delta_0}{(1+\delta_0)^2}
\).
If requiring \(h_{\gamma}(\rho)\ge 2^{-p}\), must have \(2^{-p}\le \frac{4\delta_0}{(1+\delta_0)^2}\), equivalent to \eqref{eq:xi-comoving-bit-threshold}.
Last sentence obtained by subtracting \eqref{eq:xi-comoving-bit-threshold} from \eqref{eq:xi-offcritical-endpoint-bit-asymptotic}. \qed
\end{proof}
